\chapter{Racines carrées}\label{ch:g9:sqrt}

Quelle longueur a la diagonale d'un carré de côté $1$~? Le théorème de
Pythagore répond $\sqrt 2$, un nombre dont le carré est $2$ --- et qui
se révèle ne pas être une \omterm{def:g9:fractions:fraction}{fraction}. Ce chapitre définit les \omterm{def:g9:sqrt:def}{racines
carrées}, établit les règles pour calculer avec elles, et enseigne à
simplifier des expressions comme $\sqrt{75}$.

\section{Définition et premières propriétés}

\begin{definition}[Racine carrée]\label{def:g9:sqrt:def}
Soit $a \geq 0$. La \emph{racine carrée}\index{racine carrée} de $a$,
notée $\sqrt a$, est l'unique nombre \emph{positif ou nul} dont le
carré est $a$~:
\[
\sqrt a \geq 0
\qquad\text{et}\qquad
\left(\sqrt a\right)^2 = a .
\]
Les nombres négatifs n'ont pas de racine carrée, car tout carré est
positif ou nul.
\end{definition}

\begin{example}
$\sqrt{49} = 7$, $\sqrt 0 = 0$, $\sqrt 1 = 1$,
$\sqrt{2.25} = 1.5$. Les premières \omterm{def:g9:sqrt:def}{racines carrées} à connaître par
cœur sont celles des \emph{carrés parfaits}~: $1, 4, 9, 16, 25, 36, 49,
64, 81, 100, 121, 144$.
\end{example}

\begin{omfigure}
\begin{tikzpicture}[scale=1.35]
  \draw[thick, fill=omDef!8] (0,0) rectangle (2,2);
  \draw[omThm, thick] (0,0) -- (2,2);
  \node[below, font=\small] at (1,0) {$1$};
  \node[left, font=\small] at (0,1) {$1$};
  \node[omThm, above left=-2pt, font=\small, rotate=45] at (1.08,0.92)
    {$\sqrt2$};
  \draw[thin] (1.75,0) -- (1.75,0.25) -- (2,0.25);
\end{tikzpicture}

{\small La diagonale d'un carré unité a pour longueur $\sqrt{1^2 + 1^2}
= \sqrt2 \approx 1.414$, par le théorème de Pythagore
(\cref{ch:g9:trig} le rappelle). Ce nombre est irrationnel~: ses
décimales ne se répètent jamais.}
\end{omfigure}

\begin{proposition}[Racine carrée d'un carré]\label{prop:g9:sqrt:square}
Pour tout nombre réel $x$ (positif ou non)~:
\[
\sqrt{x^2} = \abs{x} =
\begin{cases}
x & \text{si } x \geq 0,\\
-x & \text{si } x < 0 .
\end{cases}
\]
\end{proposition}

\begin{proof}
Le nombre $\abs{x}$ est positif ou nul et son carré est $x^2$ (un
nombre et son opposé ont le même carré). Étant l'unique nombre positif
ou nul de carré $x^2$, il est $\sqrt{x^2}$.
\end{proof}

\begin{example}
$\sqrt{(-5)^2} = \sqrt{25} = 5 = \abs{-5}$~: la \omterm{def:g9:sqrt:def}{racine carrée} «~oublie~»
le signe, elle ne le restaure pas. En particulier $\sqrt{x^2} = x$ est
\emph{faux} pour $x$ négatif.
\end{example}

\section{Produits et quotients de racines carrées}

\begin{theorem}[Règles de multiplication et de division]\label{thm:g9:sqrt:rules}
Pour tous $a \geq 0$ et $b \geq 0$~:
\[
\sqrt{a b} = \sqrt a \times \sqrt b,
\qquad\text{et pour } b > 0:\quad
\sqrt{\frac ab} = \frac{\sqrt a}{\sqrt b} .
\]
\end{theorem}

\begin{proof}
Le nombre $\sqrt a \times \sqrt b$ est positif ou nul (\omterm{def:g2:mult:def}{produit} de
nombres positifs ou nuls), et son carré est
\[
\left(\sqrt a \times \sqrt b\right)^2
= \left(\sqrt a\right)^2 \times \left(\sqrt b\right)^2
= a b .
\]
Par unicité du nombre positif ou nul dont le carré est $ab$, il égale
$\sqrt{ab}$. La règle du \omterm{def:g3:division:remainder}{quotient} se prouve de la même façon.
\end{proof}

\begin{remark}[Pas de telle règle pour les sommes~!]
$\sqrt{a + b}$ n'est \emph{pas} $\sqrt a + \sqrt b$ en général~:
\[
\sqrt{9 + 16} = \sqrt{25} = 5,
\qquad\text{mais}\qquad
\sqrt 9 + \sqrt{16} = 3 + 4 = 7 .
\]
\end{remark}

\begin{method}[Simplifier $\sqrt n$]\label{met:g9:sqrt:simplify}
Pour simplifier la \omterm{def:g9:sqrt:def}{racine carrée} d'un entier~:
\begin{enumerate}
  \item trouver le plus grand carré parfait divisant $n$ (décomposer
        $n$ si besoin)~;
  \item scinder avec la règle du \omterm{def:g2:mult:def}{produit} et extraire le carré~:
        $\sqrt{k^2 m} = k\sqrt m$~;
  \item vérifier qu'aucun carré parfait ne \omterm{def:g9:arith:divisor}{divise} ce qui reste sous la
        racine.
\end{enumerate}
\end{method}

\begin{example}\label{ex:g9:sqrt:simplify}
Simplifier $\sqrt{75}$~: puisque $75 = 25 \times 3$,
\[
\sqrt{75} = \sqrt{25 \times 3} = \sqrt{25} \times \sqrt 3 = 5\sqrt3 .
\]
Simplifier $\sqrt{72}$~: puisque $72 = 36 \times 2$,
$\sqrt{72} = 6\sqrt2$. Et un \omterm{def:g3:division:remainder}{quotient}~:
$\sqrt{\dfrac{49}{16}} = \dfrac{\sqrt{49}}{\sqrt{16}} = \dfrac74$.
\end{example}

\begin{example}[Additionner des racines carrées]\label{ex:g9:sqrt:add}
Les \omterm{def:g1:addition:def}{sommes} de \omterm{def:g9:sqrt:def}{racines carrées} ne se simplifient que lorsque les racines
sont \emph{semblables}. Calculer $\sqrt{12} + \sqrt{27} - \sqrt{48}$, en
simplifiant d'abord chaque terme~:
\begin{align*}
\sqrt{12} &= \sqrt{4 \times 3} = 2\sqrt3, \\
\sqrt{27} &= \sqrt{9 \times 3} = 3\sqrt3, \\
\sqrt{48} &= \sqrt{16 \times 3} = 4\sqrt3,
\end{align*}
donc la \omterm{def:g1:addition:def}{somme} est $2\sqrt3 + 3\sqrt3 - 4\sqrt3 = (2 + 3 - 4)\sqrt3 =
\sqrt3$.
\end{example}

\begin{example}[Développer avec des racines carrées]\label{ex:g9:sqrt:expand}
Les identités de l'algèbre s'appliquent aux \omterm{def:g9:sqrt:def}{racines carrées}. Développer
$\left(\sqrt5 + 2\right)^2$ avec $(a+b)^2 = a^2 + 2ab + b^2$~:
\[
\left(\sqrt5 + 2\right)^2 = 5 + 2 \times 2\sqrt5 + 4 = 9 + 4\sqrt5 .
\]
Et avec la troisième identité, $(a+b)(a-b) = a^2 - b^2$~:
\[
\left(\sqrt7 + \sqrt3\right)\left(\sqrt7 - \sqrt3\right) = 7 - 3 = 4 :
\]
le \omterm{def:g2:mult:def}{produit} de deux nombres irrationnels peut être un entier.
\end{example}

\section{L'équation \texorpdfstring{$x^2 = a$}{x2 = a}}

\begin{theorem}[Résoudre $x^2 = a$]\label{thm:g9:sqrt:equation}
\begin{enumerate}
  \item Si $a > 0$, l'\omterm{def:g8:equations:def}{équation} $x^2 = a$ a exactement deux solutions~:
        $\sqrt a$ et $-\sqrt a$~;
  \item si $a = 0$, la seule solution est $0$~;
  \item si $a < 0$, il n'y a pas de solution.
\end{enumerate}
\end{theorem}

\begin{proof}
Réécrire $x^2 = a$ comme $x^2 - \left(\sqrt a\right)^2 = 0$ lorsque
$a \geq 0$, et factoriser en \omterm{ex:g1:subtraction:difference}{différence} de carrés~:
\[
\left(x - \sqrt a\right)\left(x + \sqrt a\right) = 0 ,
\]
donc $x = \sqrt a$ ou $x = -\sqrt a$ (ils coïncident lorsque $a = 0$).
Lorsque $a < 0$ il n'y a pas de solution, car $x^2 \geq 0 > a$ pour tout
$x$.
\end{proof}

\begin{omfigure}
\begin{tikzpicture}
\begin{axis}[omaxis,
  width=7.6cm, height=5.4cm,
  xmin=-3.1, xmax=3.1, ymin=-1.6, ymax=5.4,
  xtick={-1.732,1.732}, xticklabels={$-\sqrt3$,$\sqrt3$},
  ytick={3}, yticklabels={$3$}]
  \addplot[omDef, thick, domain=-2.25:2.25, samples=90] {x^2};
  \addplot[omThm, thick, domain=-2.9:2.9] {3};
  \fill (axis cs:-1.732,3) circle (1.5pt);
  \fill (axis cs:1.732,3) circle (1.5pt);
  \draw[dashed] (axis cs:-1.732,3) -- (axis cs:-1.732,0);
  \draw[dashed] (axis cs:1.732,3) -- (axis cs:1.732,0);
\end{axis}
\end{tikzpicture}

{\small Résoudre $x^2 = 3$ sur la parabole $y = x^2$~: la droite
horizontale $y = 3$ la coupe aux deux abscisses $\pm\sqrt3$.}
\end{omfigure}

\begin{example}
$x^2 = 16$~: solutions $4$ et $-4$. \quad
$x^2 = 5$~: solutions $\sqrt5$ et $-\sqrt5$ (valeurs exactes~; $\sqrt5
\approx 2.236$). \quad
$x^2 = -9$~: pas de solution. \quad
$3x^2 = 21$~: diviser d'abord par 3, $x^2 = 7$, solutions $\pm\sqrt7$.
\end{example}

\section{Exercices}

\begin{exercise}[$\star$]\label{exo:g9:sqrt:1}
Calculer sans calculatrice~:
\[
\sqrt{64}, \qquad \sqrt{100}, \qquad \sqrt{0.09}, \qquad
\sqrt{\tfrac{1}{25}}, \qquad \left(\sqrt{13}\right)^2, \qquad
\sqrt{(-4)^2} .
\]
\end{exercise}

\begin{exercise}[$\star$]\label{exo:g9:sqrt:2}
Simplifier~:
\[
\sqrt{50}, \qquad \sqrt{45}, \qquad \sqrt{98}, \qquad \sqrt{300} .
\]
\end{exercise}

\begin{exercise}[$\star$]\label{exo:g9:sqrt:3}
Calculer et simplifier~:
\[
\sqrt2 \times \sqrt{18}, \qquad
\frac{\sqrt{75}}{\sqrt3}, \qquad
\sqrt5 \times \sqrt{20}, \qquad
\frac{\sqrt{8}}{\sqrt{50}} .
\]
\end{exercise}

\begin{exercise}[$\star$]\label{exo:g9:sqrt:4}
Résoudre~: $x^2 = 36$~; \quad $x^2 = 11$~; \quad $x^2 + 4 = 0$~; \quad
$2x^2 = 50$.
\end{exercise}

\begin{exercise}[$\star$]\label{exo:g9:sqrt:5}
Réduire en un seul terme~: $\sqrt{20} + \sqrt{45}$, puis
$3\sqrt8 - \sqrt{18} + \sqrt2$.
\end{exercise}

\begin{exercise}[$\star\star$]\label{exo:g9:sqrt:6}
Développer et simplifier~:
\[
\left(\sqrt3 + 1\right)^2, \qquad
\left(2\sqrt5 - 3\right)^2, \qquad
\left(\sqrt6 + \sqrt2\right)\left(\sqrt6 - \sqrt2\right).
\]
\end{exercise}

\begin{exercise}[$\star\star$]\label{exo:g9:sqrt:7}
Un champ carré a pour \omterm{def:g6:measure:area}{aire} $6400$ m$^2$. Quelle est la longueur de son
côté~? De sa diagonale (valeur exacte, puis arrondie au mètre le plus
proche)~?
\end{exercise}

\begin{exercise}[$\star\star$]\label{exo:g9:sqrt:8}
Montrer que $\dfrac{1}{\sqrt2} = \dfrac{\sqrt2}{2}$ (multiplier
\omterm{def:g4:fractions:def}{numérateur} et \omterm{def:g4:fractions:def}{dénominateur} par $\sqrt2$). Utiliser le même truc pour
écrire $\dfrac{6}{\sqrt3}$ et $\dfrac{10}{\sqrt5}$ sans \omterm{def:g9:sqrt:def}{racine carrée} au
\omterm{def:g4:fractions:def}{dénominateur}.
\end{exercise}

\begin{exercise}[$\star\star$]\label{exo:g9:sqrt:9}
Vrai ou faux~? Justifier par une preuve ou un contre-exemple.
\begin{enumerate}
  \item Pour tous $a, b \geq 0$~: $\sqrt{ab} = \sqrt a \sqrt b$.
  \item Pour tous $a, b \geq 0$~: $\sqrt{a+b} = \sqrt a + \sqrt b$.
  \item Pour tout $x$~: $\sqrt{x^2} = x$.
  \item $\left(3\sqrt2\right)^2 = 18$.
\end{enumerate}
\end{exercise}

\begin{exercise}[$\star\star\star$]\label{exo:g9:sqrt:10}
Soit $x = \sqrt{7 + 4\sqrt3}$. Calculer
$\left(2 + \sqrt3\right)^2$, et en déduire une expression plus simple
de $x$. Même question pour $\sqrt{9 - 4\sqrt5}$ (viser un carré de la
forme $\left(\sqrt5 - 2\right)^2$, et faire attention au signe).
\end{exercise}

\section{Problème~: le nombre qui n'est pas une fraction}

\begin{problem}[{Devoir maison --- la démonstration légendaire de
l'irrationalité de $\sqrt2$, ses nombreux cousins, et la recette
étonnamment efficace de Héron pour l'approcher}]
\label{pb:g9:sqrt:1}
La diagonale du carré unité est une longueur parfaitement réelle ---
on l'a tracée dans \cref{ex:g8:pythagoras:diagonal} --- et pourtant,
comme les pythagoriciens l'ont découvert avec horreur il y a quelque
vingt-cinq siècles, \emph{aucune \omterm{def:g9:fractions:fraction}{fraction}} ne la mesure. La légende
prétend que la découverte fut punie par noyade. Ce problème parcourt
pas à pas la démonstration immortelle (quatre questions, et elle est
à nous pour la vie), multiplie les victimes, et s'achève sur la
recette qu'ont employée les ingénieurs pendant deux mille ans pour
apprivoiser l'inapprivoisable.

\medskip
\noindent\textbf{Partie I --- La démonstration.}
\begin{enumerate}
  \item Traquons-le d'abord~: vérifier que $1.4 < \sqrt2 < 1.5$,
        puis que $1.41 < \sqrt2 < 1.42$, en élevant les bornes au
        carré. Une décimale de plus~: entre quels nombres à trois
        décimales $\sqrt2$ se situe-t-il~?
  \item Supposons maintenant --- pour aboutir à une absurdité ---
        que $\sqrt2 = \frac ab$ pour une certaine \omterm{def:g9:fractions:fraction}{fraction}
        \emph{irréductible} (\cref{met:g9:arith:simplify}). Élever
        les deux membres au carré et montrer que $a^2 = 2 b^2$.
        Quelle est la parité de $a^2$~?
  \item Les faits de parité du~\cref{pb:g8:pythagoras:1} disent que
        les \omterm{def:g3:numbers:evenodd}{nombres impairs} ont des carrés \omterm{def:g3:numbers:evenodd}{impairs}. En déduire que
        $a$ est \omterm{def:g3:numbers:evenodd}{pair}, écrire $a = 2k$, substituer --- et montrer que
        $b$ doit être \omterm{def:g3:numbers:evenodd}{pair} lui aussi.
  \item Où est la contradiction~? Conclure, et énoncer le théorème
        en entier~: \emph{$\sqrt2$ est irrationnel --- ce n'est le
        \omterm{def:g3:division:remainder}{quotient} d'aucun couple d'entiers}.
  \item La démonstration s'est appuyée une fois, discrètement, sur
        le mot «~irréductible~». Désigner l'étape exacte qui
        s'effondrerait sans lui, et expliquer pourquoi il était
        légitime de supposer la \omterm{def:g9:fractions:fraction}{fraction} irréductible dès le départ.
\end{enumerate}

\medskip
\noindent\textbf{Partie II --- Les victimes se multiplient.}
\begin{enumerate}[resume]
  \item Un second style de démonstration, par les décompositions en
        facteurs premiers (\cref{thm:g9:arith:factorization})~: dans
        $a^2 = 3b^2$, comparer la parité de l'\emph{\omterm{def:g8:powers:def}{exposant} de $3$}
        de chaque côté (\cref{pb:g9:arith:1}~: les carrés portent
        des \omterm{def:g8:powers:def}{exposants} \omterm{def:g3:numbers:evenodd}{pairs}). En conclure que $\sqrt3$ est
        irrationnel.
  \item Reprendre le même argument d'\omterm{def:g8:powers:def}{exposants} sur $a^2 = n b^2$
        pour un entier $n$ quelconque~: pour quels $n$
        produit-il une contradiction, et pour lesquels échoue-t-il~?
        Énoncer le résultat complet~: $\sqrt n$ est irrationnel
        exactement lorsque\,\dots
  \item Démontrer le lemme protecteur~: un rationnel non nul
        multiplié par un irrationnel donne un irrationnel.
        (Supposer $r x = q$ avec $r, q$ rationnels, $r \neq 0$, et
        isoler $x$.) En déduire que $2\sqrt2$ et
        $\frac{\sqrt2}{2}$ sont irrationnels.
  \item Démontrer que $1 + \sqrt2$ est irrationnel. Puis la jolie~:
        en supposant $s = \sqrt2 + \sqrt3$ rationnel, calculer
        $s^2$, isoler $\sqrt6$, et trouver la contradiction.
  \item Tempérons l'enthousiasme~: donner deux nombres irrationnels
        dont la \emph{\omterm{def:g1:addition:def}{somme}} est rationnelle, et deux dont le
        \emph{\omterm{def:g2:mult:def}{produit}} est rationnel. (L'irrationalité ne se
        conserve pas par les opérations --- chaque cas réclame sa
        propre démonstration.)
\end{enumerate}

\medskip
\noindent\textbf{Partie III --- La recette de Héron.}
Deux mille ans avant les calculatrices, Héron d'Alexandrie approchait
$\sqrt2$ ainsi~: \emph{deviner $x$~; remplacer la devinette par la
moyenne de $x$ et de $\frac2x$~; recommencer.}
\begin{enumerate}[resume]
  \item Partir de la devinette $x = 1$ et calculer les deux
        devinettes suivantes sous forme de \omterm{def:g9:fractions:fraction}{fractions} exactes.
  \item Calculer la troisième devinette, encore sous forme de
        \omterm{def:g9:fractions:fraction}{fraction} exacte, puis sa valeur décimale. Comparer avec la
        question 1~: combien de décimales de $\sqrt2$ sont déjà
        exactes~?
  \item L'idée derrière la recette~: un rectangle d'\omterm{def:g6:measure:area}{aire} $2$ dont un
        côté vaut $x$ a l'autre côté égal à $\frac2x$. Montrer que
        $\sqrt2$ se situe toujours \emph{entre} $x$ et $\frac2x$
        (examiner les deux cas $x^2 > 2$ et $x^2 < 2$)~: faire la
        moyenne des deux côtés resserre donc le rectangle vers le
        carré.
  \item Les devinettes $\frac32$, $\frac{17}{12}$,
        $\frac{577}{408}$ cachent une pépite~: calculer
        $3^2 - 2 \times 2^2$, puis $17^2 - 2 \times 12^2$, puis
        $577^2 - 2 \times 408^2$. La partie I a démontré que
        $a^2 - 2b^2 = 0$ est impossible --- de combien les \omterm{def:g9:fractions:fraction}{fractions}
        de Héron s'approchent-elles de l'impossible~?
  \item Bouquet final, en deux phrases~: la diagonale du carré unité
        existe sur le papier, aucune \omterm{def:g9:fractions:fraction}{fraction} ne la mesure, et son
        écriture décimale ne peut jamais être périodique
        (\cref{pb:g9:fractions:1}). Quelle sorte de «~nouveaux
        nombres~» la droite graduée doit-elle donc contenir --- et
        où, dans cette série, leur histoire complète est-elle
        racontée~?
\end{enumerate}
\end{problem}
